% GENERATED FILE. Edit canonical lessons, then run npm run generate:books.
% canonical-source-hash: fnv1a64:a6f0be7c706c384e
% canonical-lessons: KA-C179-taja, KA-C179-halasida, KA-C179-hasi, KA-C179-sappe, KA-C179-dubari

\chapter{Fresh, Stale, Raw, Bland, Expensive}
\label{ch:ka-fresh-stale-raw-bland-expensive}
\hlchaptermodality{179}

\begin{chapteropening}
\textbf{By the end of this chapter:} \emph{I can say fresh, stale, raw, bland and expensive.}

Complete the last lesson of chapter 179: I can say fresh, stale, raw, bland and expensive.
\end{chapteropening}

\section[\texorpdfstring{\kn{ತಾಜಾ} (tājā)}{ತಾಜಾ (tājā)}]{\kn{ತಾಜಾ} (tājā) — fresh}

\label{lesson:KA-C179-taja}



\begin{quote}
Before the new one: say the Kannada for in case, then the Kannada for by the way.
\end{quote}

\subsection*{You'll want to know: \kn{ತಾಜಾ}}
\textbf{\kn{ತಾಜಾ}} — \emph{tājā} — \textquotedblleft{}fresh\textquotedblright{}.

\begin{grammarlens}[title={What you've built}]
One more word for this chapter.
\end{grammarlens}

\subsection*{Your turn}
\begin{itemize}
\raggedright
  \item \emph{Say it:} \emph{tājā}
  \item \emph{Say it:} \emph{tājā}, once more
  \item \emph{Say it:} say \emph{tājā}
  \item \emph{Recall:} say the Kannada for in case, then the Kannada for by the way, then say \emph{tājā} again
\end{itemize}

\begin{tcolorbox}[breakable,colback=teal!4,colframe=teal!35!black,title={Before you move on}]
What does \kn{ತಾಜಾ} mean? (\textquotedblleft{}Fresh\textquotedblright{}.) Say it once more.
\end{tcolorbox}
\section[\texorpdfstring{\kn{ಹಳಸಿದ} (haḷasida)}{ಹಳಸಿದ (haḷasida)}]{\kn{ಹಳಸಿದ} (haḷasida) — stale, spoiled (food)}

\label{lesson:KA-C179-halasida}



\begin{quote}
Before the new one: say the Kannada for by the way, then the Kannada for fresh.
\end{quote}

\subsection*{You'll want to know: \kn{ಹಳಸಿದ}}
\textbf{\kn{ಹಳಸಿದ}} — \emph{haḷasida} — \textquotedblleft{}stale, spoiled (food)\textquotedblright{}.

\begin{grammarlens}[title={What you've built}]
One more word for this chapter.
\end{grammarlens}

\subsection*{Your turn}
\begin{itemize}
\raggedright
  \item \emph{Say it:} \emph{haḷasida}
  \item \emph{Say it:} \emph{haḷasida}, once more
  \item \emph{Say it:} say \emph{haḷasida}
  \item \emph{Recall:} say the Kannada for by the way, then the Kannada for fresh, then say \emph{haḷasida} again
\end{itemize}

\begin{tcolorbox}[breakable,colback=teal!4,colframe=teal!35!black,title={Before you move on}]
What does \kn{ಹಳಸಿದ} mean? (\textquotedblleft{}Stale, spoiled (food)\textquotedblright{}.) Say it once more.
\end{tcolorbox}
\section[\texorpdfstring{\kn{ಹಸಿ} (hasi)}{ಹಸಿ (hasi)}]{\kn{ಹಸಿ} (hasi) — raw, uncooked}

\label{lesson:KA-C179-hasi}



\begin{quote}
Before the new one: say the Kannada for fresh, then the Kannada for stale.
\end{quote}

\subsection*{You'll want to know: \kn{ಹಸಿ}}
\textbf{\kn{ಹಸಿ}} — \emph{hasi} — \textquotedblleft{}raw, uncooked\textquotedblright{}.

\begin{grammarlens}[title={What you've built}]
One more word for this chapter.
\end{grammarlens}

\subsection*{Your turn}
\begin{itemize}
\raggedright
  \item \emph{Say it:} \emph{hasi}
  \item \emph{Say it:} \emph{hasi}, once more
  \item \emph{Say it:} say \emph{hasi}
  \item \emph{Recall:} say the Kannada for fresh, then the Kannada for stale, then say \emph{hasi} again
\end{itemize}

\begin{tcolorbox}[breakable,colback=teal!4,colframe=teal!35!black,title={Before you move on}]
What does \kn{ಹಸಿ} mean? (\textquotedblleft{}Raw, uncooked\textquotedblright{}.) Say it once more.
\end{tcolorbox}
\section[\texorpdfstring{\kn{ಸಪ್ಪೆ} (sappe)}{ಸಪ್ಪೆ (sappe)}]{\kn{ಸಪ್ಪೆ} (sappe) — bland, tasteless}

\label{lesson:KA-C179-sappe}



\begin{quote}
Before the new one: say the Kannada for stale, then the Kannada for raw.
\end{quote}

\subsection*{You'll want to know: \kn{ಸಪ್ಪೆ}}
\textbf{\kn{ಸಪ್ಪೆ}} — \emph{sappe} — \textquotedblleft{}bland, tasteless\textquotedblright{}.

\begin{grammarlens}[title={What you've built}]
One more word for this chapter.
\end{grammarlens}

\subsection*{Your turn}
\begin{itemize}
\raggedright
  \item \emph{Say it:} \emph{sappe}
  \item \emph{Say it:} \emph{sappe}, once more
  \item \emph{Say it:} say \emph{sappe}
  \item \emph{Recall:} say the Kannada for stale, then the Kannada for raw, then say \emph{sappe} again
\end{itemize}

\begin{tcolorbox}[breakable,colback=teal!4,colframe=teal!35!black,title={Before you move on}]
What does \kn{ಸಪ್ಪೆ} mean? (\textquotedblleft{}Bland, tasteless\textquotedblright{}.) Say it once more.
\end{tcolorbox}
\section[\texorpdfstring{\kn{ದುಬಾರಿ} (dubāri)}{ದುಬಾರಿ (dubāri)}]{\kn{ದುಬಾರಿ} (dubāri) — expensive}

\label{lesson:KA-C179-dubari}



\begin{quote}
Before the new one: say the Kannada for raw, then the Kannada for bland.
\end{quote}

\subsection*{You'll want to know: \kn{ದುಬಾರಿ}}
\textbf{\kn{ದುಬಾರಿ}} — \emph{dubāri} — \textquotedblleft{}expensive\textquotedblright{}.

\begin{grammarlens}[title={What you've built}]
One more word for this chapter.
\end{grammarlens}

\subsection*{Your turn}
\begin{itemize}
\raggedright
  \item \emph{Say it:} \emph{dubāri}
  \item \emph{Say it:} \emph{dubāri}, once more
  \item \emph{Say it:} say \emph{dubāri}
  \item \emph{Recall:} say the Kannada for raw, then the Kannada for bland, then say \emph{dubāri} again
\end{itemize}

\begin{tcolorbox}[breakable,colback=teal!4,colframe=teal!35!black,title={Before you move on}]
What does \kn{ದುಬಾರಿ} mean? (\textquotedblleft{}Expensive\textquotedblright{}.) Say it once more.
\end{tcolorbox}
